\cotefile{76}
\scanpage{3}

% 1, heading
Complexes \deleted{cellulaires} \inserted{\unsure{côniques|cômiques|coniques}} et \deleted{\unsure{g}}ensembles à opérateurs

% 2, overwritten list, then above it
\[
\begin{aligned}
&\inserted{F_0\quad F_1\qquad F_n}\\
&\deleted{I_0\quad I_1\quad I_n}\qquad\text{ensembles \inserted{(finis?)}}\quad(\unsure{F|I}_\alpha\ \text{ens. des }\alpha\text{-faces})
\end{aligned}
\]

% 3, incidence
\[
\forall\ 0\leq\alpha\leq n-1,\quad I_\alpha\subset F_\alpha\times F_{\alpha+1}\qquad\text{(relation d'incidence)}
\]

% 4, axiom
\uline{Axiome 1} \quad $I_\alpha\to F_{\alpha+1}$ \deleted{et \unsure{$I_\alpha\to F_\alpha$}} est surjectif

% 4, left edge
% margin note written sideways
\margin{\unsure{faudr.t-il|faudrait-il|facultatif} ? \uline{Axiome $1'$} $I_\alpha\to F_\alpha$ surjectif.}

% 5, order
Sur $F_*=\displaystyle\coprod_{0\leq\alpha\leq n}F_\alpha$, on met la relation d'ordre suivante

% 6–9, braced alternatives
\[
\begin{aligned}
\underset{\substack{\rotatebox{90}{$\scriptstyle\ni$}\\F_\alpha}}{x}
<\underset{\substack{\rotatebox{90}{$\scriptstyle\ni$}\\F_\beta}}{y}
\quad\text{ssi}\quad
\left\{\begin{aligned}
&\text{ou bien }x=y\quad(\text{donc }\alpha=\beta)\\
&\text{ou bien }\alpha<\beta,\ \text{et }\exists\ z_\alpha=x\inserted{\in F_\alpha},z_{\alpha\unsure{-|+}1}\in F_{\alpha+1},\ldots,z_{\beta-1}\in F_{\beta-1}\\
&\hspace{2em} z_\beta\deleted{\unsure{\in}}=y\in F_\beta,\quad\text{tels que }\forall\alpha\leq i<\beta,\\
&\hspace{3em}z_i\ \text{et }z_{i+1}\ \text{soient incidents}
\end{aligned}\right.
\end{aligned}
\]

% 10–12, NB
\uline{NB} Le système $(F_\alpha)_{0\leq\alpha\leq n}$ et les $I_\alpha$ $(0\leq\alpha\leq n-1)$ se réconstituent à partir de l'ens. ordonné $F_*$, en prenant

% 13–16, successive minima and long repetition strokes
\[
\begin{aligned}
 F_0&=\text{ens des él. minimaux de }F_*\\
 F_1&=\rule{7em}{.35pt}\quad F_*-F_0\\
 &\unsure{\cdots|\vdots}\\
 F_i&=\rule{7em}{.35pt}\quad F_*-F_0-\ldots-F_{i-1}\qquad(i\leq n)
\end{aligned}
\]

% 13–16, right edge opposite the minima
% slanted
\margin{\uline{NB} Pour une \unsure{partie} de ces considérations, il est \unsure{opportun} de fixer un \unsure{$n$ de référence|ordre inverse}.}

% 17, recovered incidence
\[
 I_\alpha\subset F_\alpha\times F_{\alpha\unsure{-|+}1}\qquad
 I_\alpha=\bigl((x,y)\mid x\in F_\alpha,\ y\in F_{\alpha+1},\ x\leq y\bigr\}
\]

% 18, bracketed line
[Celà marche parceque $I_\alpha\unsure{-|\to} F_{\alpha+1}$ surjectif)

% 19–23, chain condition
Les ens. ordonnés $F_*$ ainsi obtenus sont \unsure{essentiellement} ceux satisfaisant la condition des chaines, et \unsure{.:|.à|où} \inserted{les chaines (= ss-ens. \unsure{totalemt|totalement} ordonnés de $F_{\unsure{\alpha|*}}$) sont de} \deleted{les chaines croissantes \inserted{de ens. ordonnés} d'éléments de $F_*$} longueurs \deleted{maxima} $\leq n$ ($=n$ si on exige $F_n\neq\varnothing$).

% 19–20, left beside chain condition
% slanted
\margin{\uline{vérifier}}

% 24–27, flags
\deleted{\unsure{Ai|Ax}} \uline{Drapeaux} Ce sont par définition les parties \unsure{total.t} ordonnées non vides de $I$. \uline{longueur} d'un drapeau (= son cardinal), type $\in\mathscr P[0,n]$ (c'est une partie \uline{non vide} de $[1,n]$.

% 26–27, lower right
\margin{Ens $\mathscr D(F_*)$ des drapeaux, \unsure{Ordonné|ordonnés} $\subseteq$}

% 28–29, continues at foot
par inclusion, les él. de $F_*$ sont les drapeaux minimaux, la relation $x\leq y$ de $F_*$ signifie \unsure{(équiv.}

% lower left
% archive sheet number: 2
